\section{The computation for perfect \texorpdfstring{$\mathbb F_p$}{Fp}-algebras}
\label{app:characteristic-p}
In characteristic $p$, we obtain the graded Tambara functor from that
of $\HZ$ by inverting Euler classes and extending to Witt
vectors. We also compute cyclotomic restriction.

Let $k$ be a perfect $\F_p$-algebra, $A=W(k)$, and $\varphi$ its Witt
Frobenius.
Write $\underline A$ for the constant Tambara functor.
By \cite[Proposition~2.1]{KrauseNikolaus} and \cite{Bokstedt},
$\pi_*\THH(Hk)=k[\beta]$ with $|\beta|=2$.
Thus the underlying spectrum is $p$-complete, and the cyclotomic equivalences
in Section~\ref{sec:introduction} give $\THH(Hk)\simeq X_k$ on every $C_{p^n}$.

\subsection{Geometric fixed points and the cyclotomic restriction tower}

We identify $X_k$ with the geometric $C_p$-fixed points of
$H\underline A$ by taking the limit along cyclotomic restriction.
For this comparison to preserve norms, we take the limit in genuine
commutative rings.

For a genuine commutative $C_{p^{r+1}}$-ring $Y$, the map
\begin{equation*}
\omega_Y:Y^{C_p}\longrightarrow\Phi^{C_p}Y
\end{equation*}
preserves norms by
\cite[Theorem~6.1 and Corollary~6.5]{BHNorm} and
\cite[Proposition~2.19 and Theorem~2.35]{NormTHH}.
By \cite[Proposition~1.2, Theorem~2.36 and
Theorems~4.6--4.7]{NormTHH}, the cyclotomic equivalence
$\theta:\varrho_p^*\Phi^{C_p}X_k\xrightarrow{\sim}X_k$
also preserves norms, and
\[
R_{\mathrm{cyc}}=\theta\circ\varrho_p^*(\omega_{X_k}):
\varrho_p^*(X_k^{C_p})\longrightarrow X_k.
\]
For finite cyclic groups, we omit the pullbacks along $\varrho_p$.

\begin{definition}
For $r\ge0$, take the following limit in genuine commutative
$C_{p^r}$-rings:
\[
\mathcal R_r=\holim_{j\ge0,\,R_{\mathrm{cyc}}^{C_{p^j}}}
(X_k^{C_{p^j}})|_{C_{p^r}}.
\]
Let $p_r:\mathcal R_r\to X_k|_{C_{p^r}}$ be the projection.
The underlying limit is computed in genuine spectra.
\end{definition}

We start with the following observation.

\begin{proposition}\label{prop:cp-restriction-limit}
The map
\begin{equation}\label{cp-limit-geometric}
\theta\circ\Phi^{C_p}p_{r+1}:
\Phi^{C_p}\mathcal R_{r+1}\xrightarrow{\sim}X_k|_{C_{p^r}}
\end{equation}
is an equivalence of genuine commutative rings. There is an equivalence
$\iota_r:H\underline A\xrightarrow{\sim}\mathcal R_r$ inducing
\begin{equation}\label{cp-witt-gauge}
A\longrightarrow\pi_0^{C_{p^i}}\mathcal R_r=W(k),
\qquad a\longmapsto\varphi^{1-i}(a)
\quad(0\le i\le r).
\end{equation}
\end{proposition}

\begin{proof}
Since $X_k$ is connective, \eqref{cp-limit-geometric} is an
equivalence by \cite[Proposition~5.2, the proof of Corollary~5.4,
and Corollary~5.13]{AN}.
By \cite[Theorem~6.14]{AN},
\[
\pi_s^K\mathcal R_r=
\begin{cases}W(k),&s=0,\\0,&s\ne0,\end{cases}
\qquad K\le C_{p^r},
\]
and $\res_{C_{p^i}}^{C_{p^l}}=\varphi^{l-i}$.
The residual circle acts trivially on $W(k)$.
Since restriction is invertible, restriction determines transfers
and norms. Thus \eqref{cp-witt-gauge} is an isomorphism of
Tambara functors, which induces $\iota_r$ by
\cite[Theorems~5.2--5.3]{Ullman}.
\end{proof}

The comparison in the derived category of Mackey modules is given in
\cite[Theorem~3]{Kriz}. Using $\iota_{n+1}$, we obtain a comparison of
genuine commutative $C_{p^n}$-rings.

\begin{proposition}
\label{prop:cp-normed-comparison}
For every perfect $k$ and $n\geq0$, the composite
\[
c_{n,k}:\Phi^{C_p}H\underline A_{C_{p^{n+1}}}
 \xrightarrow{\Phi^{C_p}\iota_{n+1}}\Phi^{C_p}\mathcal R_{n+1}
 \xrightarrow{\Phi^{C_p}p_{n+1}}\Phi^{C_p}X_k
 \xrightarrow{\theta}X_k|_{C_{p^n}}
\]
is an equivalence of genuine commutative $C_{p^n}$-rings, compatible
with restriction to subgroups in the homotopy category. The map on $\pi_0^{C_{p^i}}$ is
\begin{equation}\label{cp-finite-gauge}
A/p^{i+1}\longrightarrow W_{i+1}(k),\qquad
 [a]\longmapsto[\varphi^{-i}(a)].
\end{equation}
For $k=\F_p$ the source can equally be
$\Phi^{C_p}H\underline\Z$.
\end{proposition}

\begin{proof}
The equivalence follows from \eqref{cp-limit-geometric}. At the preimage
$C_{p^{i+1}}$ of $C_{p^i}$, the map on degree-zero homotopy is
$\varphi^{-i}$. Projection followed by cyclotomic
restriction is Witt truncation, so we obtain
\eqref{cp-finite-gauge}.

For $k=\F_p$, the map
$\Phi^{C_p}H\underline\Z\to\Phi^{C_p}H\underline{\Z_p}$
is a genuine equivalence by \cite[Section~1, p.~3]{ZhangHZ}.
\end{proof}

We can therefore pass from $\F_p$ to $k$ on representation-graded
homotopy groups by extending scalars from $\Z_p$ to $W(k)$.

\begin{proposition}\label{prop:cp-flat-extension}
For $V\in RO(C_{p^n})$,
\begin{equation}\label{cp-flat-extension}
 \pi_V^{C_{p^n}}X_k\cong
A\otimes_{\Z_p}\pi_V^{C_{p^n}}X_{\F_p}.
\end{equation}
This identification preserves products, restrictions and transfers.
\end{proposition}
\begin{proof}
Since $A$ is flat over $\Z_p$, write
\[
A\cong\operatorname*{colim}_{\alpha}\Z_p^{r_\alpha}
\]
as a filtered colimit of finite free modules. Thus, as genuine
$C_{p^{n+1}}$-spectra,
\[
H\underline A\simeq
\operatorname*{colim}_{\alpha}
\bigl(H\underline{\Z_p}\bigr)^{\oplus r_\alpha}.
\]
Geometric fixed points preserve this colimit. Attention that representation
spheres are compact. Proposition~\ref{prop:cp-normed-comparison}
therefore gives
\[
\begin{aligned}
\pi_V^{C_{p^n}}X_k
&\cong
\pi_V^{C_{p^n}}\Phi^{C_p}H\underline A\\
&\cong
\operatorname*{colim}_{\alpha}
\left(
\pi_V^{C_{p^n}}\Phi^{C_p}H\underline{\Z_p}
\right)^{\oplus r_\alpha}\\
&\cong
\operatorname*{colim}_{\alpha}
\left(
\Z_p^{r_\alpha}\otimes_{\Z_p}
\pi_V^{C_{p^n}}X_{\F_p}
\right)\\
&\cong
A\otimes_{\Z_p}\pi_V^{C_{p^n}}X_{\F_p}.
\end{aligned}
\]
The inverse sends $a\otimes x$ to $ax$ and is induced by
$H\underline{\Z_p}\to H\underline A$ and Proposition~\ref{prop:cp-normed-comparison},
with the $A$-action given by \eqref{cp-finite-gauge}.
It preserves products and bilinearity, restrictions by naturality, and transfers by
Frobenius reciprocity.
\end{proof}

\subsection{Derivation from $H\underline{\mathbb{Z}}$}

We use Euler localization to derive the Tambara functor for $X_k$ from $\HZ$ and then compute cyclotomic restriction. We begin by determining how many powers of the Euler
class suffice in a given representation degree.

\begin{proposition}\label{lem:cp-stabilization}
Let $q:C_{p^{n+1}}\twoheadrightarrow C_{p^n}$ be the quotient map, and
let $\lambda=\lambda_{p^{n+1}}(1)$ be the faithful oriented plane.
For $H\leq C_{p^n}$ write $\widetilde H=q^{-1}H$ and
$\lambda_H=\lambda|_{\widetilde H}$.  For every $V\in\RO(H)$,
\begin{equation}\label{cp-euler-bound}
 \pi_{q^*V-m\lambda_H}^{\widetilde H}H\underline A
       \xrightarrow{\ \sim\ }\pi_V^H X_k,
 \qquad m\geq0,\quad 2m>\dim V.
\end{equation}
The transition $m\mapsto m+1$ is multiplication by $a_{\lambda_H}$.
\end{proposition}
\begin{proof}
Since $A$ is flat over $\Z$, the argument for
\eqref{cp-flat-extension} applies over $\Z$.
For $m\ge0$ and $2m>\dim V$,
\cite[Section~1, p.~3]{ZhangHZ} and
Proposition~\ref{prop:cp-normed-comparison} give
\[
\begin{aligned}
\pi_{q^*V-m\lambda_H}^{\widetilde H}H\underline A
&\cong
A\otimes_{\Z}
\pi_{q^*V-m\lambda_H}^{\widetilde H}\HZ\\
&\xrightarrow{\sim}
A\otimes_{\Z}\pi_V^H\Phi^{C_p}\HZ\\
&\cong
\pi_V^H\Phi^{C_p}H\underline A\\
&\xrightarrow{\sim}
\pi_V^H X_k.
\end{aligned}
\]
Under the first isomorphism,
$1\otimes(a_{\lambda_H}\cdot-)$ corresponds to
$a_{\lambda_H}\cdot-$.
\end{proof}

We next describe how to take the norm of a class represented at one of
these finite stages.

\begin{proposition}\label{prop:cp-induced}
The $RO(C_{p^n})$-graded Tambara functor of $X_k$ is obtained
from the $RO(C_{p^{n+1}})$-graded Tambara functor of $\HZ$ by localization of Euler classes and extension of
coefficients to $A$.
\end{proposition}

\begin{proof}
By Propositions~\ref{prop:cp-normed-comparison}
and~\ref{prop:cp-flat-extension} and naturality, we get the Green functor structure.
It remains to determine the norms.

Use the notation of Proposition~\ref{lem:cp-stabilization}.
Let $H\le L\le C_{p^n}$ and $V\in RO(H)$.
For
\[
x\in\pi_{q^*V-m\lambda_H}^{\widetilde H}\HZ,
\qquad m\ge0,\quad 2m>\dim V,
\]
write $x/a_{\lambda_H}^m$ for its image
in $\pi_V^H X_{\F_p}$ under
Proposition~\ref{prop:cp-normed-comparison}.
Then
\[
N_H^L\!\left(\frac{x}{a_{\lambda_H}^m}\right)
=
\frac{N_{\widetilde H}^{\widetilde L}(x)}
     {a_{m\Ind_{\widetilde H}^{\widetilde L}\lambda_H}}
\in\pi_{\Ind_H^L V}^L X_{\F_p},
\]
since
\[
\begin{gathered}
N_{\widetilde H}^{\widetilde L}(a_{\lambda_H}^m)
=
a_{m\Ind_{\widetilde H}^{\widetilde L}\lambda_H},
\qquad
\bigl(\Ind_{\widetilde H}^{\widetilde L}\lambda_H\bigr)^{C_p}=0,
\\
\Ind_{\widetilde H}^{\widetilde L}(q^*V)
\cong q^*(\Ind_H^L V).
\end{gathered}
\]

We next determine the norms after extension of coefficients to $A$,
first on pure tensors and then on finite sums.
Under \eqref{cp-flat-extension}, by naturality and \eqref{cp-finite-gauge}, we have
\[
N_H^L(a\otimes x)
=
a^{[L:H]}\otimes N_H^L(x),
\qquad
a\in A,\quad x\in\pi_V^H X_{\F_p}.
\]
For $H=C_{p^h}$ and $L=C_{p^\ell}$, the degree-zero formula
in Witt coordinates is
\[
\begin{aligned}
N_H^L:W_{h+1}(k)&\longrightarrow W_{\ell+1}(k),\\
[\varphi^{-h}(a)]&\longmapsto
[\varphi^{-\ell}(a^{p^{\ell-h}})],
\qquad a\in A,
\end{aligned}
\]
which is the Witt vector norm
\cite[Theorem~4.1]{WittNorm}.

Suppose $[L:H]=p$. For $r\ge2$, translation by $L$ gives
\[
f:L/H\longrightarrow\{1,\ldots,r\},
\qquad |f(L/H)|>1
\quad\Longrightarrow\quad
\operatorname{Stab}_L(f)=H.
\]
Thus \cite[(8.34)]{ZhangHZ} gives $N_H^L$ on finite sums
from products, conjugation and $\tr_H^L$.
For natural $U\in RO(H)$,
\cite[Proposition~6.10]{ZhangHZ},
\eqref{cp-euler-bound} and \eqref{cp-flat-extension} give
\[
\pi_{U-1}^H X_k\cdot\pi_{U-1}^H X_k=0.
\]
Hence
\[
N_H^L(x+y)=N_H^L(x)+N_H^L(y),
\qquad x,y\in\pi_{U-1}^H X_k.
\]

For arbitrary inclusions, use
\[
N_H^L=N_K^L\circ N_H^K
\qquad(H\le K\le L).
\]
\end{proof}

We now compute cyclotomic restriction.
By \cite[Theorems~5.2--5.3]{Ullman}, the identity of $\underline A$ induces an equivalence
\[
d_r:(H\underline A_{C_{p^{r+1}}})^{C_p}
\xrightarrow{\sim}H\underline A_{C_{p^r}},
\qquad r\ge1.
\]
The transition maps of restriction induce
\[
(\mathcal R_{r+1})^{C_p}\xrightarrow{\sim}\mathcal R_r,
\]
which is the identity in Witt coordinates. By
\eqref{cp-witt-gauge}, the map in constant coordinates at
$C_{p^i}$ is
\[
\varphi^{i-1}\circ\operatorname{id}\circ\varphi^{-i}
=\varphi^{-1},\qquad 0\le i\le r.
\]
Naturality of $\omega$ and \cite[Lemma~5.3]{AN} identify
$R_{\mathrm{cyc}}$, under
Proposition~\ref{prop:cp-normed-comparison}, with the map induced
by $H(\varphi^{-1})d_r$.

Let $U=\rho_p(\mathbf a)$, $\mathbf a\in D_n$, and choose an even
$m>0$ with $2m>\dim U+1$. In the complex of
\cite[(5.5), Theorems~5.9, 5.12 and~6.4]{ZhangHZ},
\[
\begin{gathered}
q^*U-m\lambda
=\rho_{C_{p^{n+1}}}(a_0-m,a_0,\ldots,a_n),\\
a_0-m<-\tfrac12,\qquad \ell_0:A_0\xrightarrow{\sim}T_0.
\end{gathered}
\]
Delete $A_0,T_0$ and tensor the remaining complex with $A$ over
$\mathbb Z$. Multiplication by $a_{2\lambda}$ induces
$[q]\mapsto[q]$ on each remaining coordinate
\cite[(A.46)--(A.52)]{ZhangHZ}.
By \eqref{cp-flat-extension} and \eqref{cp-euler-bound},
the kernel and cokernel is $\pi_U^{C_{p^n}}X_k$ and
$\pi_{U-1}^{C_{p^n}}X_k$, respectively.

Assume $n\ge1$. By
\cite[Theorem~4.7 and Corollary~4.35]{ZhangHZ},
the next Euler localization deletes the first remaining
vertex and edge. The target degree is
\[
\varrho_p^*(U^{C_p})
=\rho_{C_{p^{n-1}}}(a_1,\ldots,a_n).
\]
By \cite[(A.46)--(A.52)]{ZhangHZ}, Euler localization induces
$[q]\mapsto[q]$ on each remaining coordinate.
Naturality of the connecting map gives the same projection
on cokernels. For $p=2$,
\[
R_{\mathrm{cyc}}(a_{\sigma_{2^n}})
=\begin{cases}
a_{\sigma_{2^{n-1}}},&n>1,\\
1,&n=1.
\end{cases}
\]
For $p=2$, $n=1$ and $a_1\ge0$, the terminal generator has
fractional representative $2^{a_1}$ in both complexes
\cite[(5.5), (A.52)]{ZhangHZ}.

For $V=V_0+D$, use the sphere unit $\widetilde v_D$ of
Definition~\ref{def:character-translation}.
The unit has nonequivariant degree $1$, so its image
in $H\underline{\mathbb Z}_{(p)}$ is the class of
\cite[Lemma~2.6]{ZhangHZ}.
Inflation commutes with induction along preimage subgroups,
so \cite[(2.7)]{ZhangHZ} applies.
For the maps $f_{r,j}$,
\[
\varrho_p^*(f_{r,j}^{C_p})
=\begin{cases}
\operatorname{id}_{S^0},&j=0,\\
f_{r,j-1},&j>0,
\end{cases}
\]
where $f_{r,j-1}$ on the right is formed over $C_{p^{n-1}}$.
Since $\widetilde v_{r,j}=r f_{r,j}^{-1}$,
\[
\varrho_p^*\Phi^{C_p}(\widetilde v_D)
=\left(\prod_r r^{m_{r,0}}\right)
\widetilde v_{\varrho_p^*(D^{C_p})}.
\]

We now put these results together.

\begin{theorem}\label{thm:cp-result}
For every perfect $\F_p$-algebra $k$ and $n\ge0$, the calculation for
$\HZ$ gives the groups $\pi_V^H X_k$ for all $H\le C_{p^n}$
and $V\in RO(H)$, Tambara functor, and genuine
cyclotomic restriction for $n\ge1$.
\end{theorem}
\begin{proof}
Apply Proposition~\ref{prop:cp-induced} to \cite[Theorems~1.1
and~1.4; (2.7), (5.5), (8.12)--(8.14), (8.28) and (8.31)--(8.34)]{ZhangHZ} gives the tambara functor.
We use the grading convention preceding
Proposition~\ref{prop:vertex-norms}.
For $n\ge1$ and $U=\rho_p(\mathbf a)$, $\mathbf a\in D_n$, put
$W=\varrho_p^*(U^{C_p})$. In degrees
$U\to W$ and $U-1\to W-1$, respectively, $H(\varphi^{-1})d_n$ induces
\[
 R_{\mathrm{cyc}}(x_0,\ldots,x_n)=(\varphi^{-1}x_1,\ldots,\varphi^{-1}x_n),
 \qquad R_{\mathrm{cyc}}[z_0,\ldots,z_{n-1}]=
 [\varphi^{-1}z_1,\ldots,\varphi^{-1}z_{n-1}].
\]
The second target is zero for $n=1$. Multiplicativity and the
fixed-point degrees of $\widetilde v_D$ above extends the map for degree $V_0+D$.
\end{proof}
